% =============================================================================
% Abstract (200-350 words)
% =============================================================================

\section*{Abstract}
\addcontentsline{toc}{section}{Abstract}

Community detection is a fundamental problem in network science,
with applications in social network analysis, fraud detection, and
biological systems. The practical standard remains classical
heuristics such as Louvain~\cite{blondel2008louvain} and
Leiden~\cite{traag2019leiden}, although recent work has explored
formulating community detection as a Quadratic Unconstrained Binary
Optimization (QUBO) problem and solving it on specialized hardware
such as quantum annealers and Coherent Ising Machines. These
hardware-based approaches deliver promising results but remain
inaccessible to most practitioners.

This thesis develops a software-based alternative to hardware QUBO
solvers for community detection. The starting point is the QIGNN
framework~\cite{qignn2026}~--~an unsupervised graph neural network
originally validated on two-class (single-bit-per-vertex) combinatorial problems such as
MaxCut and Maximum Independent Set. The QIGNN paradigm is extended
to the multi-class setting of community detection, combining the
QUBO formulations of modularity maximization due to Wang et
al.~\cite{wang2024quantum} with the neural-relaxation approach of
Schuetz et al.~\cite{schuetz2022combinatorial}.

The QIGNN architecture is adapted to the multi-class community
assignment task through a softmax output layer over $k$ classes and
the linearization of the QUBO problem required for handling one-hot
encoded variables. An orthogonality regularization adapted
from~\cite{bianchi2020spectral} is integrated to control the
trivial-collapse failure mode of the multi-class loss landscape, and
its behaviour is analyzed across graph structures of varying
balance. A memory-efficient training procedure is developed that
computes the QUBO loss directly from the modularity matrix and the
probability matrix, without materializing the full QUBO matrix. This
reduces the memory footprint of the multi-class loss by a factor of
$k^{2}$ in the number of communities, making training feasible on
graphs that would otherwise be out of reach.

The bipartition formulation of the proposed method converges
deterministically on the binary benchmarks considered and
reaches normalized mutual information at or above Louvain's on
the majority of them. Where the direct multi-class formulation
runs into difficulties, a hierarchical extension is proposed that
recursively applies the bipartition formulation. The hierarchical
algorithm exposes two operating modes with different trade-offs:
the fully automatic modularity-gain mode matches Louvain in
modularity within $0.02$ on every benchmark tested (up to $42$
ground-truth communities) and exceeds Louvain in NMI on four of
them; the target-$k$ mode (which requires the number of
communities as input) trades a modularity gap on imbalanced
graphs for higher NMI, matching or exceeding Louvain on ten of
thirteen benchmarks. The algorithm inherits the empirical
determinism of the bipartition formulation across both modes.
The applicability regime of each formulation is analyzed in
terms of problem size, community balance, and number of classes.

Together, these results establish QIGNN-based community detection
as a viable commodity-hardware alternative to hardware QUBO
solvers, competitive with classical heuristics across both binary
and multi-class settings.

\bigskip

\textbf{Keywords:} community detection, graph neural networks, QUBO,
modularity, combinatorial optimization, unsupervised learning.

\newpage